% !TeX spellcheck = en_US
\documentclass[french]{yLectureNote}

\title{Atomistique 2 - PS}
\subtitle{Chimie}
\author{Paulhenry Saux}
\date{\today}
\yLanguage{Français}

\professor{J.Cunny}
\usepackage{graphicx}%----pour mettre des images
\usepackage[utf8]{inputenc}%---encodage
\usepackage{geometry}%---pour modifier les tailles et mettre a4paper
%\usepackage{awesomebox}%---pour les boites d'exercices, de pbq et de croquis ---d\'esactiv\'e pour les TP de PC
\usepackage{tikz}%---pour deiffner + d\'ependance de chemfig
% \usepackage{tabularx}%---pour dimensionner automatiquement les tableaux avec variable X
\usepackage{awesomebox}%---Pour les boites info, danger et autres
\usepackage{menukeys}%---Pour deiffner les touches de Calculatrice
\usepackage{fancyhdr}%---pour les en-t\^ete personnalis\'ees
\usepackage{blindtext}%---pour les liens
\usepackage{hyperref}%---pour les liens (\`a mettre en dernier)
\usepackage{caption}%---pour la francisation de la l\'egende table vers Tableau
\usepackage{pifont}
\usepackage{array}%---pour les tableaux
\usepackage{lipsum}
\usepackage{yFlatTable}
\usepackage{multicol}
\newcommand{\Lim}[1]{\lim\limits_{\substack{#1}}\:}
\renewcommand{\vec}{\overrightarrow}
\newcommand{\N}[0]{\mathbb{N}}
\newcommand{\dd}{\mathrm{d}}
\newcommand{\norm}[1]{||\vec{#1}||}
\newcommand{\fo}{\psi(\vec{r},t)}
\newcommand{\foe}{\psi^*(\vec{r},t)}
\newcommand{\HH}{\hat{H}}
\newcommand{\hb}{\hbar}
\newcommand{\lap}{\nabla^2}
\newcommand{\lapcc}{\frac{\partial^2 }{\partial x^2}+\frac{\partial^2 }{\partial y^2}+\frac{\partial^2 }{\partial z^2}}
\begin{document}
\chapter{Rudiments de mécanique quantique}
% \section{Introduction}
% \subsection{Lois}
% \begin{theorem}[Longueur d'onde De Broglie]
%  \(\lambda = \frac{h}{mv}\) relie le caractère ondulatoire ou caractère corpusculaire.
% \end{theorem}
\section{Les postulats de la mécanique quantique}
\subsection{Fonction d'onde}
\begin{theorem}[Postulat 1]
L'état d'un système à un instant t est complètement défini par la connaissance de sa fonction d'onde notée \(\fo\).
\end{theorem}
Ainsi :
\begin{itemize}
 \item \(\foe\fo\) est la probabilité de trouver la particule au point r.\marginCritical{Une valeur en un point}
 \item \(|\fo|^2\dd V\) est la probabilité de trouver la particule sur le volume infinitésimal dV\marginCritical{Une valeur en sur une distance élémentaire}
  \item \(|\fo|^2\) est la densité de probabilité à l'instant t
\end{itemize}
\subsection{Opérateurs}
\subsubsection{Généralités}
\begin{theorem}[Postulat 2 - Opérateurs]
 À toute grandeur physique A mesurable on associe en mécanique quantique un opérateur linéaire et hermitique noté \^A. Ils prennent en entrée et sortie des fonctions d'onde.
\end{theorem}
\begin{definition}[Fonctions propres]
Une fonction propre $\psi$ de \^A est une fonction non nulle telle que l'application de \^A sur $\psi$ donne \(k \psi\)
\end{definition}
\criticalInfo{Commutativité}{Les opérateurs ne commutent pas en général. 2 opérateurs commutent si \([\hat{A},\hat{B}] = \hat{A}\cdot \hat{B} - \hat{B}\cdot \hat{A} = 0\)}
\begin{theorem}[Opérateurs commutants]
 2 opérateurs commutant admettent les m\^emes vecteurs propres
\end{theorem}
\begin{theorem}[]
Une combinaison linéaire de fonctions propres dégénérées d'un opérateur est également fonction propre de cet opérateur avec la même valeur propre
\end{theorem}
\begin{myproof}[]
Considérons une combinaison linéaire de ces fonctions propres dégénérées. On note \(k\) la valeur propre : \[\theta = c_1\psi_1+\dots+c_n\psi_n\]
Alors \begin{flalign*}
\hat{A} \theta &= \hat{A}(c_1\psi_1+\dots+c_n\psi_n)\\
&= \hat{A}c_1\psi_1+\dots+\hat{A}c_n\psi_n\\
&= kc_1\psi_1+\dots+kc_n\psi_n\\
&= k(c_1\psi_1+\dots+c_n\psi_n)\\
& = k\theta
\end{flalign*}
\end{myproof}
\subsubsection{Principe de correspondance}
\begin{theorem}[Principe de correspondance]
 Tout opérateur peut \^etre construit à partir des opérateurs position et quantité de mouvement.
\end{theorem}
\begin{definition}[Opérateur à connaitre]
\begin{itemize}
        \item \(\hat{q_i} = q_i \times\) : multiplication par la coordonnée \(q_i\)
        \item \(\hat{p_x} = -i\hbar \frac{\partial}{\partial x}\) : opérateur qt de mouv
        \item \(\hat{p} = -i\hbar \nabla\) : opérateur norme de la qt de mouv
        \item \(\hat{T} = - \frac{\hbar^2}{2m}\nabla^2 = - \frac{\hbar^2}{2m}\Delta\)
        \item \(\hat{H} = \hat{T} + \hat{V}\) avec V l'énergie potentielle de la particule.
        \item \(\hat{L}^2 = -\hb^2\Lambda\) Opérateur moment cinétique
        \item \(\hat{L_z} = -i\hb \frac{\partial}{\partial \varphi}\) Opérateur moment cinétique sur z
       \end{itemize}
\end{definition}
\subsection{Mesure physique}
\begin{theorem}[Postulat 3 - Valeurs expérimentales]
 Les valeurs de A mesurables en peuvent \^etre que les valeurs propres de l'opérateur associé \(\^A\)
\end{theorem}
\begin{proposition}[Valeur moyenne]
\label{prop1}
\[\langle A\rangle  = \frac{\int\int\int \foe\hat{A}\fo \dd V}{\int\int\int\foe \fo \dd V}\]
\end{proposition}
La dispersion sera alors
\begin{proposition}[Dispersion]
\[\Delta A = \sqrt{\langle A^2\rangle -\langle A\rangle ^2}\]
\end{proposition}
\section{Application des postulats}
\subsubsection{Inégalité d'Heisenberg}
\begin{theorem}[Inégalité d'Heinsenberg]
                      \(\Delta x\Delta p_x \geq \frac{\hb}{2}\)
                     \end{theorem}
\subsubsection{Séparation des variables}
\begin{theorem}[Séparation des variables]
 Si on peut écrire un opérateur A comme la somme de 2 opérateurs s'appliquant à 2 variables, alors le produit des vecteurs propres des deux opérateurs et le vecteur propre de l'opérateur A. La somme des valeurs propres des 2 opérateurs vaut celle de A.
\end{theorem}
\subsection{Particule sur une sphère}
\subsubsection{Analyse}
\warningInfo{Normalisation  en coordonnées sphériques}{Il faut que \[\int_0^{\infty}\int_0^{\pi}\int_0^{2\pi} \psi^*\psi r^2\sin(\theta)\dd \theta \dd \phi \dd r = 1\]}
\begin{theorem}[Fonctions propres de l'opérateur de Legendre]
\[\hat{\Lambda} Y_{lm}(\theta, \varphi) = -l(l+1)Y_{lm}(\theta, \varphi)\]
\end{theorem}
\section{Atomes hydrogénoïdes}
\subsection{Équation de Schrodinger}
L'hamiltonien est donc \[\HH = \hat{T_e} + \hat{T_N}+\hat{V_{eN}}\] On remarque cependant que l'équation ne pourra \^etre résolue car le rayon dépend de la position des 2 particules. Il faut donc simplifier l'équation.
\subsubsection{Problème à 2 corps}
En notant \(\mu = \frac{mem_N}{m_e+m_N} \simeq m_e\) la masse réduite et en se plaçant dans un repère centré\marginCheck{Cela permet à l'opérateur \(\hat{r}\) de dépendre uniquement des coordonnées de l'électron, ce qui rend l'équation solube par séparation des variables.} sur le noyau, on peut noter l'hamiltonien comme\marginCheck{Il décrit uniquement le mouvement relatif de l'électron par rapport au noyau.} \[\HH = \hat{T_{\mu}} + \hat{V_{eN}} = -\frac{\hb^2}{2\mu}\Delta \mu - \frac{1}{4\pi \epsilon_0}\frac{Z_e^2}{\hat{r}}\]
\begin{theorem}[Approximation de Born-Oppenheimer]
Du fait de leur large différence de masse, les électrons s'adaptent de façon instantanée et adiabatique (sans transfert d'énergie) à tout mouvement des noyaux : le noyau est donc fixé.
\end{theorem}
\subsubsection{Résolution}
Comme le potentiel autour du noyau présente une symétrie sphérique, on se place en coordonnées sphériques. Ainsi, en injectant l'expression de l'hamiltonien dans l'équation de Schrodinger, on obtient :
\[\HH \psi = (-\frac{\hb^2}{2m_er^2}(\frac{\partial}{\partial r}r^2\frac{\partial}{\partial r}+\Lambda)-\frac{1}{4\pi \epsilon_0}\frac{Z_e^2}{r})\psi = E\psi\]
\subsection{Fonctions propres de l'équation radiale}
\criticalInfo{Signifcation de \(\rho\)}{On pose souvent \(\rho = \frac{Zr}{a_0}\)}
\begin{theorem}[Condition de normalisation de la partie radiale]
La condition de normalisation est :
\[ \int_0^\infty |R_{1s}(r)|^2 r^2 dr = 1 \] On ajoute aussi le \(r^2\) quand on cherche à montrer l'orthogonalité.
\end{theorem}
\begin{theorem}[Densité de probabilité de présence radiale]
\[\rho_{nl}(r) = R^2_{nl}(r)r^2\]
\end{theorem}
\begin{theorem}[Valeur moyenne de la partie radiale]
\[\int_0^\infty R^2_{nl}(r)r^3 \dd r\]
\end{theorem}
\subsection{Formule de Rydberg}
L'énergie de l’électron ne dépend que de \(n\). En effet, pour un hydrogénoïde, on a

\begin{theorem}[Quantité d'énergie]
\[E_n = \frac{E_1(H)Z^2}{n^2}\]
\end{theorem}

On a également la formule de Rydberg\marginInfo{La constante \(R_H\) est une approximation car elle vient directement de la résolution de l'équation radiale, que l'on a simplifier en supposant que le noyau était fixe. Or, ce n'est pas le cas, il y a donc un décalage} :

\begin{theorem}[Formule de Rydberg]
 La longueur d'onde émise par un changement d'énergie d'un électron est reliée avec le niveau d'énergie \(n\) : \[\frac{1}{\lambda}= \mathcal{R_H}Z^2(\frac{1}{n_a^2}-\frac{1}{n^2_d})\]
\end{theorem}
\begin{theorem}[Énergie d'un état excité pour un hydrogénoïde]
 \[E_n = -\frac{13.6 Z^2}{n^2} eV\]
\end{theorem}
\subsection{Nouvelle unité}
On définit le système d'unité atomique (u.a) : On pose \(e=\hb=4\pi\epsilon_0=c=m_e=a_0=1\).

Les énergies s'écrivent alors \(E_1(H) = -\frac12 u.a.\) et \[E_n(Z) = -\frac{Z^2}{2n^2}u.a.\]
\subsection{À retenir}
\warningInfo{Nombres quantiques}{
\begin{itemize}
 \item \(n \in \mathbb{N}^+\) : nombre quantique principal
 \item \(l \in \mathbb{N}, \in [0, n-1]\) le nombre quantique angulaire
 \item \(m \in \mathbb{N}, \in [-l, l]\) par pas de 1 le nombre quantique magnétique
\end{itemize}
}
\warningInfo{Moment cinétique}{
Un électron dans une OA(l,m) a un moment cinétique de \(L = \sqrt{l(l+1)\hb}\) et une composante selon z valant \(m\hb\)}
Toutes les OA correspondant une valeur de n donnée sont dans un m\^eme couche\marginWarning{Elles ont alors le m\^eme niveau d'énergie : elles sont dites dégénérées.}Celles qui ont en plus la m\^eme valeur de l sont dans la m\^eme  sous-couche\marginWarning{m pouvant prendre 2l+1 valeurs différentes, il y a au maximum 2l+1 OA dans une sous-couche.}
\subsection{Représentation}
\subsubsection{Parties angulaires}
Les harmoniques de type s sont sphériques, les p ont des lobes autour d'axes et les d ont 4 lobes ou 2 lobes et un anneau.
\criticalInfo{Surface nodale}{Il y a l surfaces nodales pour chaque fonctions.}
\begin{theorem}[Points nodaux]
 Dans les fonctions radiales, il y a n-l-1 points nodaux. Les points nodaux en 0 ne sont pas comptés s'il y en a.
\end{theorem}
\subsection{Spin}
\begin{definition}[Spin]
Propriété intrinsèque des particules, c'est un moment cinétique intrinsèque, quantifié. On différencie le nombre quantique de spin \(s = 0.5\) et le nombre quantique magnétique de spin \(m_s = \pm 0.5\)
\end{definition}
%TODO AJouter r^2 pour l'orthogonalité ?
 \end{document}
